\chapter{Il kernel di calcolo}
\label{cap:problema}

Questo capitolo fissa tre cose che il resto della relazione dà per acquisite:
che cosa il nucleo calcola e come se ne contano le operazioni, in che ordine
conviene scandire i tre cicli, e da dove vengono i dati. Sono scelte fatte una
volta sola, prima di qualunque backend: i capitoli su MPI, CPU e CUDA le
ereditano senza rimetterle in discussione.

\section{Che cosa calcola il nucleo}
\label{sec:definizione}

Il nucleo calcola $Y \leftarrow A X$, con $A$ densa $M \times N$ e $X$
multivettore $N \times k$. Elemento per elemento:
\begin{equation}
  Y[i][c] \;=\; \sum_{j=0}^{N-1} A[i][j]\, X[j][c],
  \qquad 0 \le i < M,\ 0 \le c < k .
  \label{eq:definizione}
\end{equation}

La traccia fissa i gradi di libertà: $M$ ed $N$ sono scelte dall'utente a
runtime, $k$ deve funzionare per valore \emph{generico} e il collaudo deve
coprire almeno $k = 3,\ 6,\ 8,\ 20,\ 32$. Il progetto rispetta entrambe le
richieste, ma con una distinzione che conviene anticipare subito: sui cinque
$k$ del collaudo il nucleo di CPU usa funzioni specializzate a tempo di
compilazione, mentre ogni altro $k$ passa da un percorso generico con
prestazioni diverse (\cref{cap:kernel-cpu}).

\subsection{Conteggio delle operazioni}

Il conteggio si costruisce in due passi.

\emph{Passo 1 - un singolo elemento.} Ogni $Y[i][c]$ è una somma di $N$
prodotti: servono $N$ moltiplicazioni e $N-1$ addizioni, ma la convenzione
standard, e quella imposta dalla traccia, arrotonda a $2$ FLOP per termine.

\emph{Passo 2 - tutta la matrice.} Gli elementi di $Y$ sono $M \cdot k$.

Quindi:
\begin{equation}
  \mathrm{FLOP}(M,N,k) = 2\,M\,N\,k ,
  \qquad
  \mathrm{GFLOPS} = \frac{2\,M\,N\,k}{T \cdot 10^{9}} .
  \label{eq:flops}
\end{equation}

Su questa formula la traccia impone tre vincoli, che vale la pena enunciare
qui perché condizionano tutta la metodologia di misura:

\begin{enumerate}[nosep,left=0pt]
  \item $T$ è il tempo \textbf{medio per invocazione}, ottenuto ripetendo il
        nucleo più volte sulla stessa matrice.
  \item In $T$ entra \textbf{solo il calcolo}: preprocessamento, I/O ed
        eventuali trasferimenti da e verso la scheda CUDA restano fuori. La
        traccia consente però di misurarli e discuterli a parte, ed è quello
        che il progetto fa. Che cosa esattamente venga cronometrato, e con
        quale orologio, è trattato in \cref{sec:metriche}.
  \item $M$ ed $N$ sono le dimensioni della matrice \textbf{globale}, anche
        quando l'esecuzione è distribuita. È proprio questo che
        rende leggibile lo speedup: il lavoro totale non cambia con il numero
        di processi, cambia solo $T$.
\end{enumerate}

\section{Come vengono memorizzati i dati}
\label{sec:layout}

Tutte e tre le matrici sono \textbf{row-major}: \emph{gli elementi di una stessa
riga sono adiacenti}.

La matrice $A$ ha inoltre una propria \emph{leading dimension} \code{lda}, cioè
la distanza in elementi fra l'inizio di una riga e l'inizio della successiva
nella rappresentazione in memoria.
La leading dimension è stata introdotta per permettere \emph{padding} fra righe consecutive
e per ricevere direttamente blocchi non contigui mediante datatype MPI derivati.
In realtà, le matrici X e Y la leading dimension è \emph{sempre} fissata a $k$,
che è il numero di colonne, e quindi le righe sono compatte.
La leading dimensione è sfruttata:
\[
  A[i][j] \;=\; \code{A\_loc}[\,i \cdot \code{lda} + j\,], \qquad
  X[j][c] \;=\; \code{X\_loc}[\,j \cdot \code{ldx} + c\,], \qquad
  Y[i][c] \;=\; \code{Y\_loc}[\,i \cdot \code{ldy} + c\,].
\]
Quando \code{lda} coincide con il numero di colonne il buffer è compatto;
quando è maggiore, fra una riga e l'altra restano delle celle inutilizzate, per l'appunto il \emph{padding}.

\paragraph{Perché \code{lda} è separata da \code{n\_loc}.}
Il blocco che un processo possiede è una sottomatrice di $A$, quindi
\textbf{non è contiguo nella matrice globale}: le sue righe sono distanti $N$
elementi l'una dall'altra. La distribuzione usa per questo un datatype
derivato \code{MPI\_Type\_vector}, con stride $N$ dal lato del mittente e
stride \code{lda} dal lato del ricevente, così il blocco viene depositato
direttamente nel buffer locale senza copie intermedie (\cref{cap:mpi}).
Tenere \code{lda} come parametro indipendente aggiunge a questo la possibilità
di avere padding fra le righe locali.s

\paragraph{Perché \code{ldx} e \code{ldy} invece sono fissate a $k$.}
L'interfaccia del kernel (\cref{app:interfaccia}) le espone e ogni backend le
onora, ma il driver distribuito impone $\code{ldx} = \code{ldy} = k$, cioè
buffer compatti. La ragione è nelle collettive: il \code{MPI\_Bcast} di $X$ e
il \code{MPI\_Reduce} di $Y$ lavorano con conteggi contigui, e descrivere uno
stride richiederebbe di costruire un datatype derivato a ogni invocazione,
\emph{dentro} la regione cronometrata. $A$ invece non attraversa nessuna
collettiva, ed è per questo che è l'unica delle tre a potersi permettere il
padding.

Il fatto che le $k$ colonne di $X$ siano adiacenti in memoria non è un
dettaglio di comodo: è la proprietà su cui si regge tutta la
\cref{sec:schemi}.

\begin{ipotesi}{}
{Un backend CUDA prevede in alternativa una disposizione
\textbf{column-major} di $X$, ma come opzione di build e con effetti discussi al
momento giusto (\cref{cap:kernel-cuda}); il default, e tutto ciò che segue, è
row-major.}
\end{ipotesi}

\section{Ordine dei cicli: schema A contro schema B}
\label{sec:schemi}

La \cref{eq:definizione} contiene tre indici, $i$, $j$ e $c$.
Indipendentemente dall'ordine in cui viene
scelto di scandire gli indici, l'operazione produce lo stesso risultato,
ma il diverso ordine con il quale vengono scanditi non danno affatto le stesse prestazioni.
È il primo punto in cui il progetto si
allontana da una trascrizione letterale della definizione.

\subsection{Gli schemi candidati}

I due candidati sono:

\begin{center}
\begin{minipage}[t]{0.47\linewidth}
\textbf{Schema A} ($i \to j \to c$)
\begin{lstlisting}[language=C]
for i:
  acc[0..k) = 0
  for j:
    a = A[i][j]
    for c:
      acc[c] += a * X[j][c]
  Y[i][0..k) = acc[0..k)
\end{lstlisting}
\end{minipage}\hfill
\begin{minipage}[t]{0.47\linewidth}
\textbf{Schema B} ($i \to c \to j$)
\begin{lstlisting}[language=C]
for i:
  for c:
    s = 0
    for j:
      s += A[i][j] * X[j][c]
    Y[i][c] = s
\end{lstlisting}
\end{minipage}
\end{center}

Lo schema~B è la trascrizione letterale della definizione, ed è esattamente
l'esecuzione di $k$ prodotti matrice-vettore indipendenti, uno per colonna.
Lo schema~A scambia i due cicli interni, cambiando anche il ``ruolo'' della matrice $A$.

Si segua una singola riga $i$.

\emph{Schema B.} Il ciclo su $c$ sta fuori e quello su $j$ dentro, quindi per
ogni colonna $c$ il ciclo interno percorre l'intera riga
$A[i][0 \ldots N{-}1]$. Finita una colonna si ricomincia da capo con la
successiva, sulle stesse celle. La riga viene dunque letta $k$ volte, per un
totale di $kN$ letture di $A$ a fronte delle $N$ celle che la riga contiene.

\emph{Schema A.} Il ciclo su $j$ sta fuori e quello su $c$ dentro. Il valore
$a = A[i][j]$ viene letto una volta sola e consumato subito da tutti e $k$ gli
accumulatori, con gli aggiornamenti
$\code{acc}[c] \mathrel{+}= a \cdot X[j][c]$ per $c = 0 \ldots k-1$; poi si
passa a $j+1$ e quel valore di $A$ non serve più. La riga viene letta una volta
sola, per un totale di $N$ letture. Il riuso di $A$ vale quindi esattamente
$k$, e vale per costruzione: non dipende dalla taglia del problema né da
condizioni sul valore di $k$.

Perché il risparmio sia reale gli accumulatori devono stare in registro, e
perché il compilatore possa tenerceli deve conoscerne il numero a tempo di
compilazione: è il motivo delle specializzazioni su $k$ descritte in
\cref{sec:specializzazione}.

Si potrebbe obiettare che le riletture dello schema~B le assorbe la cache, e in
parte è così, ma solo finché la dimensione della riga è contenibile nella cache.
Una riga di $A$ occupa $N s$
byte, dove  $s$ è la dimensione in byte di ogni elemento della matrice
e $N$ è la dimensione che la traccia chiede di far crescere: il caso
interessante è proprio quello in cui la riga eccede i livelli di cache vicini
al calcolo o, sulla GPU, quello in cui la riga eccede la shared memory di un blocco. Il modello che segue
conta perciò il traffico verso la memoria principale, che è il livello che
finisce per determinare il tempo.

\subsection{La quantificazione: intensità aritmetica}

L'\emph{intensità aritmetica} $\Ai$ è il rapporto fra i FLOP eseguiti e i byte
che attraversano il livello di memoria considerato. Detta $s$ la dimensione in
byte dello scalare, e contando il solo traffico su $A$ come giustificato in
\cref{sec:traffico}:

\begin{equation}
  \Ai_{\text{A}} \;=\; \frac{2\,M N k}{M N s} \;=\; \frac{2k}{s},
  \qquad
  \Ai_{\text{B}} \;=\; \frac{2\,M N k}{k\,M N s} \;=\; \frac{2}{s} .
  \label{eq:intensita}
\end{equation}

Il numeratore è lo stesso: i FLOP da fare sono fissati dal problema. Cambia il
denominatore, perché lo schema~B legge gli $MN$ elementi di $A$ una volta per
ciascuna delle $k$ colonne.

In FP64 ($s = 8$) questo significa $\Ai_{\text{A}} = k/4$ contro
$\Ai_{\text{B}} = 1/4$ \si{\flop\per\byte}: lo schema~A guadagna esattamente
un fattore $k$, e a $k = 32$ porta l'intensità da $0{,}25$ a
\SI{8}{\flop\per\byte}. 

È questo salto che sposta il nucleo da un regime saldamente memory-bound a un
regime in cui, sulla GPU in doppia precisione, il limite può diventare il
calcolo. La discussione quantitativa, con i picchi reali della scheda, è nel
\cref{cap:roofline}.

\subsection{Un limite che non si può superare}

$\Ai_{\text{A}} = 2k/s$ è il \textbf{massimo} ottenibile: gli $MN$ elementi di
$A$ vanno letti almeno una volta e i FLOP sono fissati, quindi nessuna
riorganizzazione ulteriore può migliorare il rapporto. Tutto il lavoro di
ottimizzazione successivo, su CPU come su GPU, riguarda allora il traffico su
$X$ e la capacità di tenere i $k$ accumulatori in registro, non l'intensità
rispetto ad $A$, che è già al suo limite.

Con un'eccezione, che la traccia rende rilevante chiedendo $k$ generico. Tenere
$k$ accumulatori in registro è possibile finché $k$ è abbastanza piccolo: oltre
una certa ampiezza il nucleo di CPU blocca le colonne a gruppi di
$\code{KB} = 32$ e rilegge $A$ una volta per gruppo, cosicché in generale
\[
  \Ai_{\text{A}}(k) \;=\; \frac{2k}{s \cdot \lceil k/\code{KB} \rceil}.
\]
Per i cinque $k$ del collaudo il ciclo sui blocchi è degenere, $A$ viene letta
una volta sola e si ricade nella \cref{eq:intensita}; per $k = 64$ i passaggi
sono due e l'intensità si ferma al valore che aveva a $k = 32$. Il dettaglio,
con il microbenchmark che lo misura, è in \cref{sec:specializzazione}.

\subsection{Una proprietà che conta sull'hardware reale}

Nello schema~A entrambi i flussi di lettura sono a \textbf{passo unitario}. La
riga di $A$ è contigua per costruzione e la riga di $X$ lo è perché $X$ è
row-major con le $k$ colonne adiacenti (\cref{sec:layout}). Nessuno dei due
accessi ha stride, quindi i prefetcher hardware lavorano nelle condizioni
migliori e, sulla GPU, gli accessi delle lane di uno stesso warp sono
coalescibili. Lo schema~B, per ottenere la stessa proprietà sull'accesso a
$X$, richiederebbe $X$ column-major, cioè un layout diverso da quello che
serve altrove.

\section{Precisione}
\label{sec:precisione}

Tutto il codice usa un unico tipo \code{scalar\_t}, definito in
\file{src/common/scalar.h}, che vale \code{double} per default e \code{float}
compilando con \code{PREC=float} (il Makefile traduce la variabile nella macro
\code{-DUSE\_FLOAT}). Anche il tipo MPI segue lo scalare, attraverso la macro
\code{SCALAR\_MPI\_TYPE}, così le collettive non vanno toccate.

Il confronto FP64/FP32 è quindi un flag di compilazione e non una seconda
copia dei sorgenti: le due build eseguono esattamente lo stesso codice, che è
la sola condizione in cui il confronto misura la precisione e non due
implementazioni diverse.

\paragraph{Effetto sul modello roofline.} La scelta della precisione non è
neutra, e agisce in \textbf{due modi opposti}:

\begin{itemize}[nosep,left=0pt]
  \item dimezzando $s$ la \cref{eq:intensita} \emph{raddoppia} l'intensità
        aritmetica, da $k/4$ a $k/2$. Da solo, questo avvicinerebbe il nucleo
        al regime compute-bound;
  \item ma il punto di crossover del roofline vale $P_{\text{picco}}/B$, e su
        una GPU Turing il picco FP32 è $32$ volte il picco FP64. Il crossover
        si sposta quindi a destra di un fattore $32$.
\end{itemize}

Il secondo effetto è $16$ volte più forte del primo e prevale nettamente: in
singola precisione il nucleo finisce \textbf{saldamente memory-bound} su tutto
l'intervallo di $k$ della traccia, mentre in doppia precisione è
compute-bound già da $k = 6$. I due regimi sono qualitativamente diversi, ed è
questo il vero argomento con cui il progetto discute la precisione: la doppia
non costa il fattore $2$ del traffico, costa il fattore $32$ del picco. I
numeri sono in \cref{sec:roofline-gpu}.

\paragraph{Effetto sulla validazione.} \file{scalar.h} deriva dal tipo anche
\code{SCALAR\_EPS} e la tolleranza del confronto con l'oracolo seriale, che
non è una costante ma cresce con la lunghezza della riduzione. Il modello e la
costante misurata sono discussi in \cref{cap:validazione}.

\section{Generazione riproducibile dei dati}
\label{sec:generazione}

Le matrici di test sono generate, non lette da file. La traccia lo consente
esplicitamente, elencando la generazione fra le strategie di collaudo, e in
cambio chiede che sia \textbf{ripetibile}.

\subsection{Il generatore}

Il generatore è \code{splitmix64}, un mixer a 64 bit \textbf{senza stato}. La
differenza rispetto a un generatore pseudocasuale ordinario è tutta qui: un
generatore con stato produce una \emph{sequenza}, e per ottenere l'elemento in
posizione $p$ bisogna averlo fatto avanzare $p$ volte. Un mixer stateless va
da $(\text{seed}, \text{chiave})$ al valore in un passo solo. La chiave scelta è
l'\textbf{indice lineare globale} dell'elemento:

\begin{equation}
  A[i][j] \;=\; g\bigl(\text{seed},\ \Sigma_A,\ i \cdot N + j \bigr),
  \qquad
  X[j][c] \;=\; g\bigl(\text{seed},\ \Sigma_X,\ j \cdot k + c \bigr),
  \label{eq:gen}
\end{equation}
dove $\Sigma_A$ e $\Sigma_X$ sono due identificatori di \emph{stream} distinti, che
servono a evitare che $A$ e $X$ abbiano lo stesso pattern di valori. Il seed
viene mescolato prima con lo stream e poi con la chiave: due passaggi, perché
uno solo lascerebbe correlati i valori di chiavi vicine. Dei $64$ bit
prodotti si usano i $53$ alti, che danno un \code{double} esatto in $[0,1)$,
riscalato poi in $[-1,1)$.

Il valore di un elemento dipende dunque \emph{solo} dalla sua posizione
globale: mai dal processo che lo calcola, mai dall'ordine in cui viene
generato, mai dalla forma della griglia.

\subsection{Le tre conseguenze, tutte usate dal progetto}

\begin{enumerate}[nosep,left=0pt]
  \item ogni processo può generare \textbf{direttamente} il proprio blocco
        locale, senza che una copia globale della matrice esista mai in
        memoria. È ciò che permette di misurare taglie che su un solo nodo non
        starebbero;
  \item la stessa matrice globale si ottiene con \emph{qualunque} griglia di
        processi e con qualunque modalità di inizializzazione. Due
        configurazioni diverse sono quindi confrontabili elemento per
        elemento;
  \item l'esecuzione seriale di \cref{cap:validazione}, usata come oracolo,
        può ricostruire da sé la
        porzione che gli serve, senza ricevere dati dagli altri processi.
\end{enumerate}

\subsection{Perché i valori hanno media nulla}

I valori sono distribuiti in $[-1, 1)$ a media nulla, e questa non è una scelta
estetica. Ogni elemento di $Y$ è una riduzione su $N$ termini, e l'errore di
arrotondamento di una riduzione dipende da come si compongono i segni: con
termini a media nulla gli arrotondamenti hanno segno indipendente e l'errore
cresce come $\varepsilon\sqrt{N}$, non come il pessimistico $\varepsilon N$
del caso in cui tutti i termini abbiano lo stesso segno. È questa proprietà
dei dati che rende possibile una tolleranza di validazione stretta invece di
una soglia fissa e generosa (\cref{cap:validazione}).

Una nota utile per il confronto fra precisioni: il generatore calcola sempre in
\code{double} e converte in \code{scalar\_t} solo alla fine. Le build FP64 e
FP32 lavorano quindi sugli stessi dati a meno del solo arrotondamento della
conversione, ed è la condizione che rende confrontabili i loro tempi.